\subsection{理想的乘法}

设 A 是一个环，a 和 b 是 A 的双边理想。形如 \(x_{1}y_{1} + \cdots + x_{n}y_{n}\) （其中 \(n \geqslant 0\) ， \(x_{i} \in a\) 和 \(y_{i} \in b\) ，对 \(1 \leqslant i \leqslant n\) ）的元素组成的集合显然是 A 的一个双边理想，记为 ab，称为双边理想 a 和 b 的乘积。在此乘法下，A 的双边理想的集合是一个幺半群，其单位元为双边理想 A。若 a, b, c 是 A 的双边理想，则 \(a(b + c) = ab + ac\) , \((b + c)a = ba + ca\) 。若 A 是交换的，则理想的乘法是交换的。

\(\mathsf{ab}\subset \mathsf{aA}\subset \mathsf{a}\) 和 \(\mathsf{ab}\subset \mathsf{Ab}\subset \mathsf{b}\) ，因此

\[
a b \subset a \cap b.\tag{21}
\]

命题 6. 设 \(a, b_{1}, \ldots, b_{n}\) 是 A 的双边理想。若 \(A = a + b_{i}\) 对所有 \(i\) 成立，则 \(A = a + b_{1}b_{2} \ldots b_{n} = a + (b_{1} \cap b_{2} \cap \cdots \cap b_{n})\) 。

由 (21) 可知只需证明 \(\mathbf{A} = \mathfrak{a} + \mathfrak{b}_1\mathfrak{b}_2\ldots \mathfrak{b}_n\) 。由归纳法，只需考虑 \(n = 2\) 的情形。由假设，存在 \(a, a' \in \mathfrak{a}, b_1 \in \mathfrak{b}_1, b_2 \in \mathfrak{b}_2\) ，使得 \(1 = a + b_1 = a' + b_2\) 。则

\[
1 = a ^ {\prime} + (a + b _ {1}) b _ {2} = \left(a ^ {\prime} + a b _ {2}\right) + b _ {1} b _ {2} \in \mathfrak {a} + \mathfrak {b} _ {1} \mathfrak {b} _ {2},
\]

从而 \(\mathbf{A} = \mathfrak{a} + \mathfrak{b}_1\mathfrak{b}_2\) 。

命题 7. 设 \(\mathfrak{b}_1, \ldots, \mathfrak{b}_n\) 是 A 的双边理想，满足 \(\mathfrak{b}_i + \mathfrak{b}_j = \mathbf{A}\) （对 \(i \neq j\) ）。则 \(\mathfrak{b}_1 \cap \mathfrak{b}_2 \cap \dots \cap \mathfrak{b}_n = \sum_{\sigma \in \mathfrak{E}_n} \mathfrak{b}_{\sigma(1)} \mathfrak{b}_{\sigma(2)} \ldots \mathfrak{b}_{\sigma(n)}\) 。特别地，若 A 是交换的，则 \(\mathfrak{b}_1 \cap \mathfrak{b}_2 \cap \dots \cap \mathfrak{b}_n = \mathfrak{b}_1 \mathfrak{b}_2 \ldots \mathfrak{b}_n\) （参见习题 2）。

首先假设 \(n = 2\) 。存在 \(a_1 \in \mathfrak{b}_1\) ， \(a_2 \in \mathfrak{b}_2\) ，使得 \(a_1 + a_2 = 1\) 。若 \(x \in \mathfrak{b}_1 \cap \mathfrak{b}_2\) ，则 \(x = x(a_1 + a_2) = xa_1 + xa_2 \in \mathfrak{b}_2\mathfrak{b}_1 + \mathfrak{b}_1\mathfrak{b}_2\) 。因此 \(\mathfrak{b}_1 \cap \mathfrak{b}_2 = \mathfrak{b}_1\mathfrak{b}_2 + \mathfrak{b}_2\mathfrak{b}_1\) 。

现在假设命题的等式对所有整数 \(<n\) 均已成立。由命题 6， \(\mathfrak{b}_{n} + (\mathfrak{b}_{1}\mathfrak{b}_{2}\ldots \mathfrak{b}_{n - 1}) = A\) ，因此

\[
\begin{array}{l} \mathfrak {b} _ {1} \cap \mathfrak {b} _ {2} \cap \dots \cap \mathfrak {b} _ {n} = (\mathfrak {b} _ {1} \cap \mathfrak {b} _ {2} \cap \dots \cap \mathfrak {b} _ {n - 1}) \mathfrak {b} _ {n} + \mathfrak {b} _ {n} (\mathfrak {b} _ {1} \cap \mathfrak {b} _ {2} \cap \dots \cap \mathfrak {b} _ {n - 1}) \\ = \Big (\sum_ {\sigma \in \mathfrak {E} _ {n - 1}} \mathfrak {b} _ {\sigma (1)} \mathfrak {b} _ {\sigma (2)} \dots \mathfrak {b} _ {\sigma (n - 1)} \Big) \mathfrak {b} _ {n} \\ \qquad + \mathfrak {b} _ {n} \Big (\sum_ {\sigma \in \mathfrak {E} _ {n - 1}} \mathfrak {b} _ {\sigma (1)} \mathfrak {b} _ {\sigma (2)} \dots \mathfrak {b} _ {\sigma (n - 1)} \Big) \\ \subset \sum_ {\tau \in \mathfrak {E} _ {n}} \mathfrak {b} _ {\tau (1)} \mathfrak {b} _ {\tau (2)} \dots \mathfrak {b} _ {\tau (n)} \subset \mathfrak {b} _ {1} \cap \mathfrak {b} _ {2} \cap \dots \cap \mathfrak {b} _ {n}. \end{array}
\]

